% Copy this abstract into metadata.yaml (or metadata.tex on Overleaf) so that
% the manuscript and submission metadata stay synchronized.
\begin{abstract}
Deep learning for radiologic image analysis is a rapidly growing field in
biomedical research and is likely to become a standard practice in modern
medicine. We used the publicly available NIH ChestX-ray14 dataset, containing
112,120 frontal-view chest X-rays labeled for 14 thoracic pathologies in order to
reproduce the 14-pathology extension of CheXNet and built a modified model,
DACNet. The reproduced CheXNet model achieved a mean test AUROC of 0.8125,
compared with 0.8414 calculated from the 14 per-pathology AUROCs reported in the
original article. The modified model, DACNet achieved a mean test AUROC of 0.8474. As a secondary
analysis, per-pathology decision thresholds selected only on the validation set
produced test macro F1 scores of 0.2813 for the CheXNet reproduction and 0.3414
for DACNet. The original pneumonia-versus-radiologist F1 experiment could not
be reproduced because its expert annotations for the 420-image test set are not
publicly available.
\end{abstract}

\section{Introduction}

Chest X-ray classification is an important task in current medical image analysis. Recently, 
deep learning models have been trained to detect various thoracic diseases from
radiographic scans. CheXNet, a landmark study in this field,
introduced a 121-layer DenseNet convolutional neural network that reportedly
outperformed radiologists specifically in pneumonia detection\cite{rajpurkar2017}.
The researchers used the NIH ChestX-ray14 dataset, which is publicly available and has over
100,000 frontal-view chest X-rays labeled with up to 14 disease
classes\cite{nihxray2017}. The success of CheXNet has inspired further
research because it demonstrated a significant step toward using artificial
intelligence to assist in clinical diagnosis, especially in regions where
access to licensed radiologists is limited\cite{hwang2023}. In the midst of a
reproducibility crisis in academia, independent researchers must reproduce
groundbreaking studies like this in order to protect the accuracy of future
research\cite{noorden2015}.

Researchers in this project set out to replicate the original CheXNet model as closely
as possible and evaluate AUC-ROC scores across all 14 disease classes. All of the code
for our models and our evaluation pipeline is publicly available in our
\href{https://github.com/dstrick17/DacNet}{GitHub repository}. While our
primary goal was to reproduce the original CheXNet study, we also recognize the
importance of its successor, CheXNeXt, which validated a similar model against
board-certified radiologists on a curated internal dataset\cite{rajpurkar2018}.
Although that expert-labeled test set is not currently publicly available, its findings
emphasize the clinical relevance of these models and continue to reinforce the need for
reproducible benchmarking on public datasets such as NIH ChestX-ray14.

Our first contribution is our reproduction of the 14-pathology CheXNet
experiment using ImageNet-pretrained DenseNet-121, a patient-wise split, and
the training details reported in the original article. Our second contribution
is DACNet, an extension that retains DenseNet-121 while incorporating focal
loss, the AdamW optimizer, color jitter, and validation-selected per-pathology
F1 thresholds. Unlike the original CheXNet study, which reported an F1 score
for pneumonia using an expert-labeled subset, our secondary analysis computes
per-pathology F1 scores across all 14 diseases using only the public NIH labels.

\section{Related Work}

Deep learning methods have proven to achieve substantial accuracy in image
classification. In 2012, AlexNet demonstrated the effectiveness of
convolutional neural networks on large-scale image classification
tasks\cite{krizhevsky2012}. Our main source of inspiration is ``CheXNet:
Radiologist-Level Pneumonia Detection on Chest X-Rays with Deep Learning.'' In
addition to CheXNet, the authors later introduced CheXNeXt, a related model
evaluated directly against radiologists on 14 thoracic diseases. Their results
showed radiologist-level performance on most conditions. However, the curated
validation set used in CheXNeXt is not publicly available, making full
replication difficult.

Despite substantial advancements in medical imaging, there is still room for
improvement\cite{shen2017}. Other studies have reported the efficacy of deep
learning in chest X-ray classification\cite{kufel2023}, ophthalmology
\cite{kermany2018}, and fracture detection in musculoskeletal radiographs
\cite{jones2020}.

\section{Scope of the Reproduction}

The original CheXNet article reports two related but also distinct experiments. The
primary experiment is a binary pneumonia detection. This primary experiment uses a single-output
DenseNet-121, a class-weighted binary cross-entropy loss, and an expert-labeled
test set of 420 images to compare the model with four radiologists. The second
experiment extends the network to all 14 pathologies by replacing the final
layer with 14 outputs and optimizing the sum of unweighted binary
cross-entropy losses. It uses a 70\%/10\%/20\% patient-wise split and reports
per-pathology AUROC.

This study reproduces the second 14-pathology experiment. The distinction is
important: the weighted loss and reported pneumonia F1 comparison belong to
the first binary pneumonia experiment. We cannot reproduce that comparison
because the expert annotations for its 420-image test set were not released.
We therefore treat the 14 per-pathology AUROCs as the primary reproduction
targets. The F1 scores reported here are a secondary analysis using public NIH
labels and are not presented as a reproduction of the radiologist comparison.

We first contacted the email address listed in the original paper, which
was associated with Stanford University. After receiving no response, we
learned that the author had moved to Harvard Medical School and contacted his
current laboratory in October 2026 to request access to the expert annotations.
We have not yet received a response.

\section{Methods}

\subsection{Dataset and patient split}

All data used in this study came from the NIH ChestX-ray14 dataset, the same
dataset used to build the CheXNet algorithm. The dataset is available through
the \href{https://www.kaggle.com/datasets/nih-chest-xrays/data}{Kaggle NIH
Chest X-rays page}\cite{nihxray2017} and the
\href{https://nihcc.app.box.com/v/ChestXray-NIHCC}{NIH distribution}
\cite{nihpress2017}. It contains 112,120 frontal-view chest X-rays from 30,805
patients. Each image is labeled with ``No Finding'' or with one or more of 14
pathologies. The labels were extracted automatically from radiology reports
and therefore contain label noise.

One major challenge is the extreme class imbalance within the dataset. While
there are only 14 possible classifications, any given image can exhibit any
combination of these diseases. This creates a theoretical space of $2^{14}$,
or 16,384 possible label combinations. Exploratory data analysis found 836
observed combinations. Approximately 54\% of the images are labeled ``No
Finding,'' and many individual combinations occur rarely. This imbalance
presents a substantial challenge for threshold-dependent metrics such as F1.

We sorted the unique patient identifiers and created a deterministic split
with seed 42. The training, validation, and held-out test sets contain 70\%,
10\%, and 20\% of patients, respectively, with no patient overlap. This
produced 21,563 training patients (78,614 images), 3,081 validation patients
(11,212 images), and 6,161 test patients (22,294 images). The exact assignments
are stored in \texttt{patient\_splits.csv} under
\texttt{reproducibility/replicate\_chexnet} and were shared by both models.

The original article states the same 70\%/10\%/20\% patient-level proportions
for the 14-pathology experiment but does not publish the exact patient
assignments. Our seeded split therefore reproduces the stated procedure rather
than the authors' exact partition.

\subsection{CheXNet reproduction}

We initialized a DenseNet-121 with ImageNet weights and replaced its classifier
with a fully connected layer producing 14 logits. Images were converted to
RGB, resized to $224\times224$, converted to tensors, and normalized using the
ImageNet channel means and standard deviations. Random horizontal flipping was
applied during training. These choices follow the preprocessing described in
the original CheXNet article.

Following the original 14-pathology extension, the loss was the sum of 14
unweighted binary cross-entropy losses, averaged over the minibatch. We used
Adam with $\beta_1=0.9$, $\beta_2=0.999$, an initial learning rate of 0.001,
and a batch size of 16. The learning rate was reduced by a factor of 10 whenever
validation loss failed to improve after an epoch. Training was configured for
at most 20 epochs with early stopping after five epochs without improvement,
and the checkpoint with the lowest validation loss was selected.

\subsection{DACNet extension}

DACNet retains the ImageNet-pretrained DenseNet-121 backbone and 14-output
classifier. It replaces binary cross-entropy with focal loss
($\gamma=2$, $\alpha=1$) to address class imbalance and uses AdamW with a
learning rate of 0.00005 and weight decay of $10^{-5}$. A
\texttt{ReduceLROnPlateau} scheduler reduces the learning rate by a factor of
10 after one epoch without improvement in validation loss. The model was
trained for at most 25 epochs with a batch size of 8 and early stopping after
five epochs without improvement in mean validation AUROC. The checkpoint with
the highest mean validation AUROC was selected.

DACNet training augmentation consisted of a random resized crop to
$224\times224$, random horizontal flipping, and color jitter with brightness
and contrast factors of 0.1. Validation and test images were resized to 256
pixels and center-cropped to $224\times224$. All images used ImageNet
normalization. DACNet used the exact patient assignments generated for the
CheXNet reproduction.

\subsection{Evaluation}

The primary outcome for the reproduction study was the area under the receiver operating 
characteristic curve (AUROC), which was calculated independently for each pathology 
on the held-out test set. Mean AUROC was calculated as the unweighted mean of the 14 pathology-level 
values. We compared these results with the values reported in the original CheXNet article.

The F1 scores require a decision threshold and are reported as a secondary analysis. For
each pathology, we selected the threshold that maximized F1 on the validation
set and then applied that threshold unchanged to the held-out test set. We
report the unweighted macro average across all 14 pathology-level F1 scores.
For transparency, the CheXNet reproduction also records its test F1 at a fixed
threshold of 0.5. No threshold was selected using the test labels.

The two models use differently scaled loss functions, so their numerical test
losses are not directly comparable. We therefore compare the models using
AUROC and F1 rather than raw loss values.

\section{Results}

\subsection{Reproduction of the 14-pathology AUROC results}

The reproduced CheXNet model achieved a mean test AUROC of 0.8125. The mean of
the 14 per-pathology values reported in the original article is 0.8414. The
reproduction exceeded the corresponding published AUROC for two pathologies,
Consolidation and Effusion, and was lower for the other 12 pathologies.
Table~\ref{tab:test_auc_scores} reports the complete comparison.

Our reproduction fell short of the published mean AUROC by 0.029 and was
lower for 12 of the 14 pathologies. As a result, we consider it as a partial reproduction
rather than an exact reproduction. The overall pattern was recovered, but the
numerical values were not. We cannot identify the cause of the gap from a
single run. However, there are several plausible explanations that follow from differences
between our setup and the original. First, the original patient-level
assignment was not published, so our seeded 70\%/10\%/20\% split
necessarily places different patients in the test set; because many
pathologies are rare, small changes in which positive cases fall in the
test set can shift per-pathology AUROC noticeably. Additionally, neural network
training is stochastic, and we completed only one run per model, so some
of the difference may be ordinary seed-to-seed variation that we did not
measure. Third, our software libraries and hardware differ from those
available in 2017, which can alter pretrained-weight handling, numerical
behavior, and data-loading details. Finally, the NIH labels are extracted
automatically from reports and are noisy, which may make results more
sensitive to the particular split. The gap was largest for Hernia,
Emphysema, Mass, Nodule, and Pneumothorax, but we did not test whether any
of these explanations accounts for it.

\begin{table*}[ht]
\centering
\small
\begin{tabular}{lrrr}
\hline
\textbf{Pathology} & \textbf{Original CheXNet} & \textbf{Reproduction} & \textbf{DACNet} \\
\hline
Atelectasis        & 0.809 & 0.790 & \textbf{0.828} \\
Cardiomegaly       & \textbf{0.925} & 0.917 & 0.913 \\
Consolidation      & 0.790 & 0.800 & \textbf{0.824} \\
Edema              & 0.888 & 0.883 & \textbf{0.895} \\
Effusion           & 0.864 & 0.874 & \textbf{0.883} \\
Emphysema          & \textbf{0.937} & 0.877 & 0.920 \\
Fibrosis           & 0.805 & 0.782 & \textbf{0.827} \\
Hernia             & 0.916 & 0.847 & \textbf{0.951} \\
Infiltration       & \textbf{0.735} & 0.695 & 0.712 \\
Mass               & \textbf{0.868} & 0.816 & 0.862 \\
Nodule             & 0.780 & 0.722 & \textbf{0.794} \\
Pleural thickening & 0.806 & 0.781 & \textbf{0.808} \\
Pneumonia          & \textbf{0.768} & 0.739 & 0.763 \\
Pneumothorax       & \textbf{0.889} & 0.850 & 0.884 \\
\hline
Mean               & 0.841 & 0.812 & \textbf{0.847} \\
\hline
\end{tabular}
\caption{Held-out test AUROC by pathology. Original values are those reported
by Rajpurkar et al. The bold value in each row is the largest of the three.}
\label{tab:test_auc_scores}
\end{table*}

\subsection{DACNet extension}

DACNet achieved a mean test AUROC of 0.8474 and exceeded the published CheXNet
AUROC for 8 of the 14 pathologies. The largest numerical gains over our
reproduction occurred for Hernia, Fibrosis, Nodule, Emphysema, and Mass.
Because this study contains one completed training run for each model and no
confidence intervals or repeated-seed analysis, these differences are
descriptive and should not be interpreted as evidence of statistical
significance.

Since our CheXNet reproduction fell short of the published results, we
wanted to build a model that improved on it and so future researchers
could continue to iterate on. We chose modifications aimed at the class
imbalance described above. Most images in the NIH dataset are labeled "no finding" 
and several pathologies are rare, so we expected the unweighted binary cross-entropy
loss to be dominated by easy negative examples. Focal loss down-weights
well-classified examples and was a natural candidate for this setting. We
chose AdamW because its decoupled weight decay offers a simple form of
regularization on a network with a large number of parameters, and color
jitter was added as a mild augmentation to reduce overfitting to the
intensity variations between scans. These choices were exploratory and
were not tuned individually. By documenting them openly, we hope to give
future open-source researchers a concrete starting point for further
improvements.

\subsection{Secondary F1 analysis}

At the fixed threshold of 0.5, the CheXNet reproduction had a test macro F1 of
0.1271. Selecting one threshold per pathology on the validation set increased
its held-out test macro F1 to 0.2813. Using the same validation-only threshold
selection procedure, DACNet achieved a test macro F1 of 0.3414. The
pathology-level values are shown in Table~\ref{tab:f1_by_disease}.

These F1 values do not reproduce the F1 score reported for the original binary
pneumonia model. The original value was calculated on expert annotations for a
different model and a 420-image test set that is not publicly available. The
present values instead describe thresholded performance of the 14-output
models against the public NIH labels.

\begin{table*}[ht]
\centering
\small
\begin{tabular}{lrr}
\hline
\textbf{Pathology} & \textbf{CheXNet reproduction} & \textbf{DACNet} \\
\hline
Atelectasis        & 0.365 & \textbf{0.413} \\
Cardiomegaly       & \textbf{0.401} & 0.394 \\
Consolidation      & 0.228 & \textbf{0.245} \\
Edema              & 0.231 & \textbf{0.240} \\
Effusion           & 0.516 & \textbf{0.531} \\
Emphysema          & 0.363 & \textbf{0.472} \\
Fibrosis           & 0.131 & \textbf{0.165} \\
Hernia             & 0.103 & \textbf{0.444} \\
Infiltration       & 0.404 & \textbf{0.424} \\
Mass               & 0.320 & \textbf{0.397} \\
Nodule             & 0.233 & \textbf{0.321} \\
Pleural thickening & 0.200 & \textbf{0.225} \\
Pneumonia          & 0.077 & \textbf{0.086} \\
Pneumothorax       & 0.367 & \textbf{0.422} \\
\hline
Macro average      & 0.281 & \textbf{0.341} \\
\hline
\end{tabular}
\caption{Held-out test F1 using pathology-specific thresholds selected only on
the validation set.}
\label{tab:f1_by_disease}
\end{table*}

\section{Reproducibility and Limitations}

The repository provides the training scripts, dependency files, exact patient
assignments, machine-readable result files, and commands for local and Kaggle
execution. The NIH dataset can be attached directly to a Kaggle notebook, so a
reviewer does not need to download all images to a personal computer. A short
smoke test checks dataset discovery, model initialization, training,
evaluation, and result output; it is not intended to reproduce the final
metrics. Full runs require GPU computation.

The exact patient assignment and trained weights from the original CheXNet
article were not available. Although our split follows the reported
70\%/10\%/20\% patient-level procedure, it is not necessarily identical to the
original split. The public NIH labels were extracted automatically from reports
and may contain errors. The expert annotations required for the original
pneumonia-versus-radiologist comparison are also unavailable. Finally, our
comparison is based on a single training run per model, so it does not quantify
variation across random seeds.


\section{Discussion}

Simply reproducing CheXNet's 14-output structure did not yield the exact
published AUROC values. Nevertheless, the reproduction recovered the overall
pattern of performance: Cardiomegaly, Effusion, Emphysema, and Pneumothorax
were among the stronger classes, while Infiltration, Nodule, and Pneumonia were
more difficult. The lower mean AUROC demonstrates why implementation details,
data partitions, and saved artifacts are important when evaluating a
replication.

DACNet produced a higher mean AUROC and validation-thresholded macro F1 than
our CheXNet reproduction on the same patient partition. However, the study does
not isolate the effect of each modification. DACNet should therefore be
interpreted as an exploratory extension rather than as part of the CheXNet
reproduction itself.

We are unsure exactly why our AUROC values differ from the original CheXNet results.
It is possible that future reproduction attempts may achieve results closer
to those originally reported, or they may produce results more consistent
with our reproduction. The F1 scores cannot be reliably compared since
the original study used an expert-annotated subset for the pneumonia-versus-radiologist comparison, 
which is not available to us. However, the F1 scores we calculated were still significantly lower
and not at expert radiologist levels.

\section{Conclusion}

This project aimed to reproduce and extend CheXNet, a landmark deep learning
model in medical imaging. We reproduced the public 14-pathology experiment and
obtained a mean held-out test AUROC of 0.8125, below the 0.8414 mean calculated
from the original article's reported values. We could not reproduce the
separate pneumonia-versus-radiologist F1 experiment because its expert labels
are not public. DACNet, an exploratory extension using focal loss, AdamW,
color jitter, and validation-selected thresholds, achieved a mean test AUROC
of 0.8474 and a macro F1 of 0.3414 on the same patient split.

In addition to the technical contributions, this project emphasizes the importance
of reproducibility in machine learning research. Our primary motivation was not only to
evaluate the claims of a high-impact study, but also to contribute openly to
the scientific community. Our codebase, exact patient split, evaluation tools,
and machine-readable results have been made publicly available in hopes of
promoting transparency and enabling further development by others.

\section{Funding and Competing Interests}

This research did not receive any specific grant from funding agencies in the
public, commercial, or not-for-profit sectors. The authors declare that study
sponsors had no role in the study design; in the collection, analysis, or
interpretation of data; in the writing of the manuscript; or in the decision
to submit the manuscript for publication. The authors have no conflicts of
interest to declare.

\section{Generative AI Usage Statement}

Some language in this manuscript was improved using generative AI tools to
clarify sentence structure and enhance readability. All AI-generated
suggestions were critically reviewed and edited by the authors to ensure
accuracy and appropriateness.
